% Academica --- headnote
% academica, 45 BC, Cumanum / Rome (dedicated to M. Terentius Varro)

\textit{The \textit{Academica} belongs to the extraordinary outpouring
of philosophical writing with which Cicero filled the year 45 BC ---
the months after the death of his daughter Tullia, and under the
dictatorship of Caesar, when the Forum was closed to him and he turned
the whole force of his mind to giving Rome a philosophy in its own
language. Its subject is the theory of knowledge: whether anything can
be known at all, what the criterion of truth might be, and how the
schools descended from Plato divided over the question. The work has
come down to us only in ruins. Cicero wrote it first in two books, the
\textit{Catulus} and the \textit{Lucullus}; then, at Varro's wish and
through Atticus's brokering, he recast it on a larger scale in four,
with Varro himself as the chief speaker. Of the first edition only the
second book, the \textit{Lucullus}, survives; of the second only the
opening book printed here, the dialogue men call ``the Varro,'' which
breaks off in mid-sentence at the portrait of Carneades.}

\textit{The setting is Cicero's villa near Cumae, where Varro --- the
most learned of the Romans, whose antiquarian labours had, as Cicero
tells him, restored to his countrymen the knowledge of who and where
they were --- has come to call, with Atticus making the third. Pressed
to say why, for all his learning, he has never written philosophy in
Latin, Varro answers that the educated will read the Greeks and the
unlearned will read nothing; Cicero presses back, and the frame becomes
the occasion for the doctrine. Varro expounds the teaching of the Old
Academy and the Peripatetics as a single inheritance --- the threefold
division of philosophy into ethics, natural philosophy, and logic, the
account of matter and quality, the criterion lodged in the trained
senses --- the synthesis Cicero had learned at Athens from Antiochus of
Ascalon. Then Cicero takes up the sceptical reply of the New Academy:
that Arcesilas, confronting the same darkness that had driven Socrates
and the older natural philosophers to confess their ignorance, denied
that anything could be grasped, and held that the wise man's part is to
withhold assent and rein in the recklessness that approves what it has
not understood.}

\textit{Much of the work's interest is in watching Cicero forge a Latin
philosophical vocabulary as he goes: coining \textit{qualitas} for the
Greek \textit{poiotes}, weighing \textit{species} against the Platonic
\textit{idea}, trying \textit{visum} and \textit{comprehendibile} for
the Stoic terms of perception and proposing them, half-apologetically,
to be ``worn smooth'' by use until they sit easily in the language.
What survives is a torso, and a difficult one --- the technical heart
of the epistemology lies partly in the lost \textit{Lucullus} --- but
even the fragment shows the design entire: a Roman, writing in Roman
words, setting the two great answers of the Greek schools side by side
and declining, as a man of the New Academy must, to pronounce between
them with more certainty than the matter allows.}
